% Options for packages loaded elsewhere
\PassOptionsToPackage{unicode}{hyperref}
\PassOptionsToPackage{hyphens}{url}
\documentclass[
  12pt,
]{article}
\usepackage{xcolor}
\usepackage[margin=1in]{geometry}
\usepackage{amsmath,amssymb}
\setcounter{secnumdepth}{-\maxdimen} % remove section numbering
\usepackage{iftex}
\ifPDFTeX
  \usepackage[T1]{fontenc}
  \usepackage[utf8]{inputenc}
  \usepackage{textcomp} % provide euro and other symbols
\else % if luatex or xetex
  \usepackage{unicode-math} % this also loads fontspec
  \defaultfontfeatures{Scale=MatchLowercase}
  \defaultfontfeatures[\rmfamily]{Ligatures=TeX,Scale=1}
\fi
\usepackage{lmodern}
\ifPDFTeX\else
  % xetex/luatex font selection
  \setmainfont[]{Arial}
\fi
% Use upquote if available, for straight quotes in verbatim environments
\IfFileExists{upquote.sty}{\usepackage{upquote}}{}
\IfFileExists{microtype.sty}{% use microtype if available
  \usepackage[]{microtype}
  \UseMicrotypeSet[protrusion]{basicmath} % disable protrusion for tt fonts
}{}
\usepackage{setspace}
\makeatletter
\@ifundefined{KOMAClassName}{% if non-KOMA class
  \IfFileExists{parskip.sty}{%
    \usepackage{parskip}
  }{% else
    \setlength{\parindent}{0pt}
    \setlength{\parskip}{6pt plus 2pt minus 1pt}}
}{% if KOMA class
  \KOMAoptions{parskip=half}}
\makeatother
\usepackage{color}
\usepackage{fancyvrb}
\newcommand{\VerbBar}{|}
\newcommand{\VERB}{\Verb[commandchars=\\\{\}]}
\DefineVerbatimEnvironment{Highlighting}{Verbatim}{commandchars=\\\{\}}
% Add ',fontsize=\small' for more characters per line
\usepackage{framed}
\definecolor{shadecolor}{RGB}{248,248,248}
\newenvironment{Shaded}{\begin{snugshade}}{\end{snugshade}}
\newcommand{\AlertTok}[1]{\textcolor[rgb]{0.94,0.16,0.16}{#1}}
\newcommand{\AnnotationTok}[1]{\textcolor[rgb]{0.56,0.35,0.01}{\textbf{\textit{#1}}}}
\newcommand{\AttributeTok}[1]{\textcolor[rgb]{0.13,0.29,0.53}{#1}}
\newcommand{\BaseNTok}[1]{\textcolor[rgb]{0.00,0.00,0.81}{#1}}
\newcommand{\BuiltInTok}[1]{#1}
\newcommand{\CharTok}[1]{\textcolor[rgb]{0.31,0.60,0.02}{#1}}
\newcommand{\CommentTok}[1]{\textcolor[rgb]{0.56,0.35,0.01}{\textit{#1}}}
\newcommand{\CommentVarTok}[1]{\textcolor[rgb]{0.56,0.35,0.01}{\textbf{\textit{#1}}}}
\newcommand{\ConstantTok}[1]{\textcolor[rgb]{0.56,0.35,0.01}{#1}}
\newcommand{\ControlFlowTok}[1]{\textcolor[rgb]{0.13,0.29,0.53}{\textbf{#1}}}
\newcommand{\DataTypeTok}[1]{\textcolor[rgb]{0.13,0.29,0.53}{#1}}
\newcommand{\DecValTok}[1]{\textcolor[rgb]{0.00,0.00,0.81}{#1}}
\newcommand{\DocumentationTok}[1]{\textcolor[rgb]{0.56,0.35,0.01}{\textbf{\textit{#1}}}}
\newcommand{\ErrorTok}[1]{\textcolor[rgb]{0.64,0.00,0.00}{\textbf{#1}}}
\newcommand{\ExtensionTok}[1]{#1}
\newcommand{\FloatTok}[1]{\textcolor[rgb]{0.00,0.00,0.81}{#1}}
\newcommand{\FunctionTok}[1]{\textcolor[rgb]{0.13,0.29,0.53}{\textbf{#1}}}
\newcommand{\ImportTok}[1]{#1}
\newcommand{\InformationTok}[1]{\textcolor[rgb]{0.56,0.35,0.01}{\textbf{\textit{#1}}}}
\newcommand{\KeywordTok}[1]{\textcolor[rgb]{0.13,0.29,0.53}{\textbf{#1}}}
\newcommand{\NormalTok}[1]{#1}
\newcommand{\OperatorTok}[1]{\textcolor[rgb]{0.81,0.36,0.00}{\textbf{#1}}}
\newcommand{\OtherTok}[1]{\textcolor[rgb]{0.56,0.35,0.01}{#1}}
\newcommand{\PreprocessorTok}[1]{\textcolor[rgb]{0.56,0.35,0.01}{\textit{#1}}}
\newcommand{\RegionMarkerTok}[1]{#1}
\newcommand{\SpecialCharTok}[1]{\textcolor[rgb]{0.81,0.36,0.00}{\textbf{#1}}}
\newcommand{\SpecialStringTok}[1]{\textcolor[rgb]{0.31,0.60,0.02}{#1}}
\newcommand{\StringTok}[1]{\textcolor[rgb]{0.31,0.60,0.02}{#1}}
\newcommand{\VariableTok}[1]{\textcolor[rgb]{0.00,0.00,0.00}{#1}}
\newcommand{\VerbatimStringTok}[1]{\textcolor[rgb]{0.31,0.60,0.02}{#1}}
\newcommand{\WarningTok}[1]{\textcolor[rgb]{0.56,0.35,0.01}{\textbf{\textit{#1}}}}
\usepackage{longtable,booktabs,array}
\newcounter{none} % for unnumbered tables
\usepackage{calc} % for calculating minipage widths
% Correct order of tables after \paragraph or \subparagraph
\usepackage{etoolbox}
\makeatletter
\patchcmd\longtable{\par}{\if@noskipsec\mbox{}\fi\par}{}{}
\makeatother
% Allow footnotes in longtable head/foot
\IfFileExists{footnotehyper.sty}{\usepackage{footnotehyper}}{\usepackage{footnote}}
\makesavenoteenv{longtable}
\usepackage{graphicx}
\makeatletter
\newsavebox\pandoc@box
\newcommand*\pandocbounded[1]{% scales image to fit in text height/width
  \sbox\pandoc@box{#1}%
  \Gscale@div\@tempa{\textheight}{\dimexpr\ht\pandoc@box+\dp\pandoc@box\relax}%
  \Gscale@div\@tempb{\linewidth}{\wd\pandoc@box}%
  \ifdim\@tempb\p@<\@tempa\p@\let\@tempa\@tempb\fi% select the smaller of both
  \ifdim\@tempa\p@<\p@\scalebox{\@tempa}{\usebox\pandoc@box}%
  \else\usebox{\pandoc@box}%
  \fi%
}
% Set default figure placement to htbp
\def\fps@figure{htbp}
\makeatother
% definitions for citeproc citations
\NewDocumentCommand\citeproctext{}{}
\NewDocumentCommand\citeproc{mm}{%
  \begingroup\def\citeproctext{#2}\cite{#1}\endgroup}
\makeatletter
 % allow citations to break across lines
 \let\@cite@ofmt\@firstofone
 % avoid brackets around text for \cite:
 \def\@biblabel#1{}
 \def\@cite#1#2{{#1\if@tempswa , #2\fi}}
\makeatother
\newlength{\cslhangindent}
\setlength{\cslhangindent}{1.5em}
\newlength{\csllabelwidth}
\setlength{\csllabelwidth}{3em}
\newenvironment{CSLReferences}[2] % #1 hanging-indent, #2 entry-spacing
 {\begin{list}{}{%
  \setlength{\itemindent}{0pt}
  \setlength{\leftmargin}{0pt}
  \setlength{\parsep}{0pt}
  % turn on hanging indent if param 1 is 1
  \ifodd #1
   \setlength{\leftmargin}{\cslhangindent}
   \setlength{\itemindent}{-1\cslhangindent}
  \fi
  % set entry spacing
  \setlength{\itemsep}{#2\baselineskip}}}
 {\end{list}}
\usepackage{calc}
\newcommand{\CSLBlock}[1]{\hfill\break\parbox[t]{\linewidth}{\strut\ignorespaces#1\strut}}
\newcommand{\CSLLeftMargin}[1]{\parbox[t]{\csllabelwidth}{\strut#1\strut}}
\newcommand{\CSLRightInline}[1]{\parbox[t]{\linewidth - \csllabelwidth}{\strut#1\strut}}
\newcommand{\CSLIndent}[1]{\hspace{\cslhangindent}#1}
\setlength{\emergencystretch}{3em} % prevent overfull lines
\providecommand{\tightlist}{%
  \setlength{\itemsep}{0pt}\setlength{\parskip}{0pt}}
\usepackage{bookmark}
\IfFileExists{xurl.sty}{\usepackage{xurl}}{} % add URL line breaks if available
\urlstyle{same}
\hypersetup{
  pdftitle={Fitting ODEs in Stan (with Python)},
  pdfauthor={Tara Hameed (adapted by Craig Fouts)},
  hidelinks,
  pdfcreator={LaTeX via pandoc}}

\title{Fitting ODEs in Stan (with Python)}
\author{Tara Hameed (adapted by Craig Fouts)}
\date{2026-03-13}

\begin{document}
\maketitle

\setstretch{2}
\section{Introduction}\label{introduction}

This script is intended to be an example of fitting Ordinary
Differential Equation (ODE) models using CmdStanPy (Python interface to
CmdStan), where Stan is a probabilistic programming language for
performing Bayesian inference and CmdStan is its command line interface.
Often when modelling we want to be able to quantify the uncertainty of
our estimates. Bayesian inference allows us to do this in a principled
manner for ODE models. All technical details mentioned here closely
follow the \href{https://mc-stan.org/docs/stan-users-guide/}{Stan} (Stan
Development Team 2025) user manual and Bayesian Data Analysis textbook
(Gelman et al. 2013). For a more detailed explanation on fitting ODEs in
Stan, please refer to the
\href{https://mc-stan.org/cmdstanpy/}{CmdStanPy} and
\href{https://mc-stan.org/docs/stan-users-guide/}{Stan} (Stan
Development Team 2025) documentation and user manual.

Before I go through the example, let's import the relevant libraries and
set a seed.

The example in this script is about fitting a simple logistic ODE to
some pre-generated fake data. Below, I outline what the simple ODE is
({[}The ODE{]}), what the fake data looks like (\hyperref[the-data]{The
Data}), how to fit the ODE to the data
(\hyperref[fitting-the-model]{Fitting the Model} and
\hyperref[checking-model-fit]{Checking Model Fit}), and finally showcase
exploring estimated parameters ({[}Estimated Parameters{]}).

\section{The Model}\label{the-model}

A simple logistic growth model (Verhulst 1838) is a popular
one-dimensional ODE that is often used for modelling population growth:

\[
\frac{df}{dt} = \beta f \left( 1 - \frac{f}{K} \right)
\]

The mean population, \(f\), grows over time, \(t\), at a growth rate of
\(\beta\). Growth slows as the population reaches its theoretical
limiting level, \(K\), often called the carrying capacity.

\section{The Data}\label{the-data}

There is some pre-generated fake data in the \texttt{logistic\_data.Rda}
file, which is in the same repository as this script. The data was
actually generated using the logistic model and Stan, but, for the
purposes of this example, we can imagine that the data, \(y\), is the
concentration of a population growing over time, \(t\), that is recorded
in an experiment where we assume the data collection method introduces
some observational noise. The aim is to fit the logistic ODE to this
collected data, estimate important parameter values, like the growth
rate \(\beta\), and also quantify the uncertainty of the estimates.

\subsection{Visualisation}\label{visualisation}

First, let's read in the data from the \texttt{logistic\_data.Rda} file:

\begin{Shaded}
\begin{Highlighting}[]
\NormalTok{logistic\_data }\OperatorTok{=}\NormalTok{ pyreadr.read\_r(}\StringTok{\textquotesingle{}logistic\_data.Rda\textquotesingle{}}\NormalTok{)[}\VariableTok{None}\NormalTok{]}
\NormalTok{Markdown(logistic\_data.head().to\_markdown(index}\OperatorTok{=}\VariableTok{False}\NormalTok{, numalign}\OperatorTok{=}\StringTok{\textquotesingle{}center\textquotesingle{}}\NormalTok{, floatfmt}\OperatorTok{=}\StringTok{\textquotesingle{}.6f\textquotesingle{}}\NormalTok{))}
\end{Highlighting}
\end{Shaded}

{\def\LTcaptype{none} % do not increment counter
\begin{longtable}[]{@{}cc@{}}
\toprule\noalign{}
y & time \\
\midrule\noalign{}
\endhead
\bottomrule\noalign{}
\endlastfoot
0.795570 & 0.000000 \\
0.174560 & 0.500000 \\
0.774322 & 1.000000 \\
3.083851 & 1.500000 \\
2.163227 & 2.000000 \\
\end{longtable}
}

and then plot the data to explore what it looks like:

\begin{Shaded}
\begin{Highlighting}[]
\NormalTok{plt.grid(alpha}\OperatorTok{=}\FloatTok{.5}\NormalTok{)}
\NormalTok{plt.scatter(logistic\_data[}\StringTok{\textquotesingle{}time\textquotesingle{}}\NormalTok{], logistic\_data[}\StringTok{\textquotesingle{}y\textquotesingle{}}\NormalTok{], s}\OperatorTok{=}\DecValTok{100}\NormalTok{, zorder}\OperatorTok{=}\DecValTok{2}\NormalTok{)}
\NormalTok{plt.xticks((}\FloatTok{0.}\NormalTok{, }\FloatTok{2.5}\NormalTok{, }\FloatTok{5.}\NormalTok{, }\FloatTok{7.5}\NormalTok{, }\FloatTok{10.}\NormalTok{))}\OperatorTok{;}
\NormalTok{plt.yticks((}\FloatTok{0.}\NormalTok{, }\FloatTok{2.5}\NormalTok{, }\FloatTok{5.}\NormalTok{, }\FloatTok{7.5}\NormalTok{, }\FloatTok{10.}\NormalTok{))}\OperatorTok{;}
\NormalTok{plt.xlabel(}\StringTok{\textquotesingle{}Time\textquotesingle{}}\NormalTok{)}
\NormalTok{plt.ylabel(}\StringTok{\textquotesingle{}y\textquotesingle{}}\NormalTok{)}
\NormalTok{plt.tight\_layout()}
\end{Highlighting}
\end{Shaded}

\pandocbounded{\includegraphics[keepaspectratio]{01_fit_ode_files/figure-latex/unnamed-chunk-2-173.pdf}}

\subsection{Pre-processing}\label{pre-processing}

Next, the dataframe needs to be re-formatted to be compatible with Stan
because Stan expects the data to be in a particular format:

\begin{Shaded}
\begin{Highlighting}[]
\CommentTok{\# Only want to solve the ODE at unique times in Stan to reduce runtime}
\NormalTok{logistic\_unique\_times }\OperatorTok{=}\NormalTok{ np.sort(logistic\_data[}\StringTok{\textquotesingle{}time\textquotesingle{}}\NormalTok{].unique())}

\CommentTok{\# Create a dictionary of data that Stan needs to define the model using the same}
\CommentTok{\# names as those in the "logistic.stan" file}
\NormalTok{stan\_logistic\_data }\OperatorTok{=}\NormalTok{ \{}
    \StringTok{\textquotesingle{}N\textquotesingle{}}\NormalTok{: }\BuiltInTok{len}\NormalTok{(logistic\_data),}
    \StringTok{\textquotesingle{}T\textquotesingle{}}\NormalTok{: }\BuiltInTok{len}\NormalTok{(logistic\_unique\_times),}
    \StringTok{\textquotesingle{}y\textquotesingle{}}\NormalTok{: logistic\_data[}\StringTok{\textquotesingle{}y\textquotesingle{}}\NormalTok{],}
    \StringTok{\textquotesingle{}time\textquotesingle{}}\NormalTok{: logistic\_unique\_times,}
    \StringTok{\textquotesingle{}time\_idx\textquotesingle{}}\NormalTok{: [np.where(logistic\_unique\_times }\OperatorTok{==}\NormalTok{ t)[}\DecValTok{0}\NormalTok{].item() }\OperatorTok{+} \DecValTok{1} \ControlFlowTok{for}\NormalTok{ t }\KeywordTok{in}\NormalTok{ logistic\_data[}\StringTok{\textquotesingle{}time\textquotesingle{}}\NormalTok{]],}
    \StringTok{\textquotesingle{}include\_likelihood\textquotesingle{}}\NormalTok{: }\VariableTok{True}
\NormalTok{\}}
\end{Highlighting}
\end{Shaded}

Although not relevant in this simple example, in the Stan model if there
are multiple observations at one time point (like multiple replicates)
the ODE is only solved once and so the data list is set up to reflect
this.

\section{Fitting the Model}\label{fitting-the-model}

\subsection{Python code}\label{python-code}

To fit the logistic ODE using Stan and estimate parameters, a Stan file
needs to be created. One has already been created and is in this
repository under the name \texttt{logistic.stan}. A copy of the file is
under the ``Stan code'' tab so that you can easily refer to the relevant
sections of the Stan model.

In the Stan file, the desired likelihood and priors can be specified.
One of the simplest \textbf{likelihoods} to use is a Gaussian, where we
assume that our data \(y_t\) at time \(t\) is normally distributed
around our ODE solution, \(f(t)\), with some noise that has scale
\(\sigma\):

\[
y_t \sim \text{N}(f(t), \sigma);
\]

In the Stan file (\texttt{logistic.stan}, see ``Stan code'' tab), the
above is written in the \texttt{model} block as follows, where
\texttt{mu} is the ODE solution:

\begin{verbatim}
y ~ normal(mu[time_idx, 1], sigma)
\end{verbatim}

We can also assume the following \textbf{priors}:

\begin{align*}
    k &\sim \text{N}^+(10, 1) \\
    \beta &\sim \text{N}^+(0, 1) \\
    \sigma &\sim \text{N}^+(0, 1) \\
    f_0 &\sim \text{N}^+(0, 1)
\end{align*}

by writing the following in \texttt{logistic.stan}'s \texttt{model}
block:

\begin{verbatim}
K ~ normal(10, 1);
beta ~ std_normal();
sigma ~ std_normal();
y0 ~ std_normal();
\end{verbatim}

These priors are chosen arbitrarily for this example, but prior choice
is important! Further information on choosing priors can be found in
(Gelman et al. 2020) or (Gelman et al. 2013).

Finally, to fit the model we only need to call the following command:

\begin{Shaded}
\begin{Highlighting}[]
\NormalTok{fit\_logistic }\OperatorTok{=}\NormalTok{ CmdStanModel(stan\_file}\OperatorTok{=}\StringTok{"logistic.stan"}\NormalTok{).sample(stan\_logistic\_data, chains}\OperatorTok{=}\DecValTok{4}\NormalTok{, seed}\OperatorTok{=}\NormalTok{seed)}
\end{Highlighting}
\end{Shaded}

Here, I've hidden the output as it is long but Stan will output a
progress bar of sampling from the posterior with the time taken for
sampling given at the end of the output. Any warnings about convergence
will also be given here.

\subsection{Stan code}\label{stan-code}

In this tab is a copy of the \texttt{logistic.stan} file used to fit the
logistic ODE in the ``Python code'' tab. It has six different
``blocks'': (\textbf{1}) \texttt{functions}, where the ODE is specified,
(\textbf{2}) \texttt{data}, where the input data is specified (notice
how the names in our dictionary \texttt{stan\_logistic\_data} match this
block), (\textbf{3}) \texttt{parameters} for the free parameters in the
model, (\textbf{4}) \texttt{transformed\ parameters}, where we solve the
ODE, (\textbf{5}) \texttt{model} that specifies the prior and
likelihood, and, finally, (\textbf{6}) \texttt{generated\ quantities},
where we simulate ODE solutions and model samples from the fitted model.

\begin{verbatim}
functions {
    vector logistic(real t, vector y, real beta, real K) {
        vector[1] dydt;

        dydt[1] = beta*y[1]*(1 - y[1]/K);

        return dydt;
    }
}

data {
    int<lower=0> N;  // Number of data points
    int<lower=0> T;  // Number of time points
    vector[N] y;  // Input data
    array[T] real time;
    array[N] int<lower=1, upper=T> time_idx;  // Index of time vector in data
    int<lower=0, upper=1> include_likelihood;
}

parameters {
    real<lower=0> K;  // Carrying capacity
    real<lower=0> beta;  // Growth rate
    real<lower=0> y0;  // Initial condition
    real<lower=0> sigma;  // Noise scale
}

transformed parameters {
    // Solve ODE
    array[T - 1] vector[1] mu_hat = ode_rk45(logistic, to_vector({y0}), time[1], time[2:T], beta, K);
    array[T] vector[1] mu;

    mu[1, 1] = y0;
    mu[2:T, 1] = mu_hat[, 1];
}

model {
    K ~ normal(10, 1);
    beta ~ std_normal();
    sigma ~ std_normal();
    y0 ~ std_normal();

    if (include_likelihood) y ~ normal(mu[time_idx, 1], sigma);
}

generated quantities {
   array[N] real f_reps = mu[time_idx, 1];  // Posterior samples of ODE solution
   array[N] real y_reps = normal_rng(mu[time_idx, 1], sigma);  // Posterior predictive
}
\end{verbatim}

\section{Checking Model Fit}\label{checking-model-fit}

In the above \texttt{sample} method call, the keyword argument
\texttt{chains} was set to be \texttt{chains=4} to specify that we want
to run 4 chains independently. This is mainly done so that we can assess
if the chains have converged to (what we assume to be) the true
posterior distribution.

In reality, we cannot know if the chains have converged to the actual
true posterior as we do not know what the true posterior is. Instead,
since of non-convergence are monitored. One way to do this is by
plotting the posterior samples over iterations for each chain (``trace
plot'') to make sure that the chains are well \textbf{mixed} and
\textbf{stationary}. The chains should all converge to the same
posterior (be well \emph{mixed}) and the chains should actually converge
to a distribution and stay there (be \emph{stationary}) (Gelman et al.
2013). If the chains are not mixed or not stationary (like if the
chains' sampling paths appear to be completely distinct or drift, for
example), these are signs of non-convergence and samples taken from
these chains should not be used further.

\begin{Shaded}
\begin{Highlighting}[]
\NormalTok{az.plot\_trace(fit\_logistic, var\_names}\OperatorTok{=}\NormalTok{(}\StringTok{\textquotesingle{}K\textquotesingle{}}\NormalTok{, }\StringTok{\textquotesingle{}beta\textquotesingle{}}\NormalTok{, }\StringTok{\textquotesingle{}y0\textquotesingle{}}\NormalTok{, }\StringTok{\textquotesingle{}sigma\textquotesingle{}}\NormalTok{), compact}\OperatorTok{=}\VariableTok{False}\NormalTok{, legend}\OperatorTok{=}\VariableTok{True}\NormalTok{)}\OperatorTok{;}
\NormalTok{plt.tight\_layout()}
\end{Highlighting}
\end{Shaded}

\pandocbounded{\includegraphics[keepaspectratio]{01_fit_ode_files/figure-latex/unnamed-chunk-5-175.pdf}}

The traceplots look OK. However, visual assessment can be subjective so
convergence is also assessed using the Gelman-Rubin metric
\(\widehat{R}\) (Vehtari et al. 2021), which can be thought of as a
numerical value summarising the above. Usually, a value of
\(\widehat{R}\) \textgreater{} 1.01 is used as a heuristic to diagnose
lack of convergence.

\begin{Shaded}
\begin{Highlighting}[]
\NormalTok{Markdown(fit\_logistic.summary().head()[}\DecValTok{1}\NormalTok{:].to\_markdown(numalign}\OperatorTok{=}\StringTok{\textquotesingle{}center\textquotesingle{}}\NormalTok{, floatfmt}\OperatorTok{=}\StringTok{\textquotesingle{}.4f\textquotesingle{}}\NormalTok{))}
\end{Highlighting}
\end{Shaded}

{\def\LTcaptype{none} % do not increment counter
\begin{longtable}[]{@{}
  >{\raggedright\arraybackslash}p{(\linewidth - 22\tabcolsep) * \real{0.0609}}
  >{\centering\arraybackslash}p{(\linewidth - 22\tabcolsep) * \real{0.0783}}
  >{\centering\arraybackslash}p{(\linewidth - 22\tabcolsep) * \real{0.0696}}
  >{\centering\arraybackslash}p{(\linewidth - 22\tabcolsep) * \real{0.0870}}
  >{\centering\arraybackslash}p{(\linewidth - 22\tabcolsep) * \real{0.0696}}
  >{\centering\arraybackslash}p{(\linewidth - 22\tabcolsep) * \real{0.0696}}
  >{\centering\arraybackslash}p{(\linewidth - 22\tabcolsep) * \real{0.0783}}
  >{\centering\arraybackslash}p{(\linewidth - 22\tabcolsep) * \real{0.0783}}
  >{\centering\arraybackslash}p{(\linewidth - 22\tabcolsep) * \real{0.1043}}
  >{\centering\arraybackslash}p{(\linewidth - 22\tabcolsep) * \real{0.1043}}
  >{\centering\arraybackslash}p{(\linewidth - 22\tabcolsep) * \real{0.1217}}
  >{\centering\arraybackslash}p{(\linewidth - 22\tabcolsep) * \real{0.0783}}@{}}
\toprule\noalign{}
\begin{minipage}[b]{\linewidth}\raggedright
\end{minipage} & \begin{minipage}[b]{\linewidth}\centering
Mean
\end{minipage} & \begin{minipage}[b]{\linewidth}\centering
MCSE
\end{minipage} & \begin{minipage}[b]{\linewidth}\centering
StdDev
\end{minipage} & \begin{minipage}[b]{\linewidth}\centering
MAD
\end{minipage} & \begin{minipage}[b]{\linewidth}\centering
5\%
\end{minipage} & \begin{minipage}[b]{\linewidth}\centering
50\%
\end{minipage} & \begin{minipage}[b]{\linewidth}\centering
95\%
\end{minipage} & \begin{minipage}[b]{\linewidth}\centering
ESS\_bulk
\end{minipage} & \begin{minipage}[b]{\linewidth}\centering
ESS\_tail
\end{minipage} & \begin{minipage}[b]{\linewidth}\centering
ESS\_bulk/s
\end{minipage} & \begin{minipage}[b]{\linewidth}\centering
R\_hat
\end{minipage} \\
\midrule\noalign{}
\endhead
\bottomrule\noalign{}
\endlastfoot
K & 10.0740 & 0.0187 & 0.9852 & 0.9857 & 8.4813 & 10.0676 & 11.7159 &
2800.6100 & 2563.5700 & 3252.7400 & 1.0025 \\
beta & 0.3363 & 0.0015 & 0.0602 & 0.0568 & 0.2485 & 0.3309 & 0.4406 &
1726.8500 & 1928.6200 & 2005.6300 & 1.0013 \\
y0 & 0.9293 & 0.0069 & 0.2948 & 0.2931 & 0.4808 & 0.9128 & 1.4508 &
1810.9300 & 1845.1600 & 2103.2800 & 1.0022 \\
sigma & 1.3199 & 0.0045 & 0.2131 & 0.2061 & 1.0226 & 1.2931 & 1.7009 &
2287.8600 & 2043.6700 & 2657.2100 & 1.0005 \\
\end{longtable}
}

The \(\widehat{R}\) values listed in the last column of the above table
look to be \textless{} 1.01!

\section{Plotting Model Fit}\label{plotting-model-fit}

During fitting, it was assumed that the data, \(y_t\), are normally
distributed around the ODE solution, \(f(t)\), with some noise that has
scale \(\sigma\) at time \(t\). Since there are no signs of
non-convergence so far, we can plot posterior samples of \(f(t)\) to
explore if they are similar to the mean data and thus check that our
fitted model can at least visually capture the general trends of the
data:

\begin{Shaded}
\begin{Highlighting}[]
\NormalTok{keys }\OperatorTok{=}\NormalTok{ [}\SpecialStringTok{f\textquotesingle{}f\_reps[}\SpecialCharTok{\{}\NormalTok{i}\SpecialCharTok{\}}\SpecialStringTok{]\textquotesingle{}} \ControlFlowTok{for}\NormalTok{ i }\KeywordTok{in} \BuiltInTok{range}\NormalTok{(}\DecValTok{1}\NormalTok{, }\DecValTok{22}\NormalTok{)]}
\NormalTok{sample\_draws }\OperatorTok{=}\NormalTok{ np.random.randint(}\DecValTok{0}\NormalTok{, }\DecValTok{4000}\NormalTok{, }\DecValTok{100}\NormalTok{)}
\NormalTok{f }\OperatorTok{=}\NormalTok{ fit\_logistic.draws\_pd()[keys].loc[sample\_draws].T.to\_numpy()}
\NormalTok{plt.plot(np.arange(}\BuiltInTok{len}\NormalTok{(f))}\OperatorTok{/}\DecValTok{2}\NormalTok{, f, alpha}\OperatorTok{=}\FloatTok{.5}\NormalTok{, linewidth}\OperatorTok{=}\FloatTok{.5}\NormalTok{, zorder}\OperatorTok{=}\DecValTok{2}\NormalTok{)}
\NormalTok{plt.grid(alpha}\OperatorTok{=}\FloatTok{.5}\NormalTok{)}
\NormalTok{plt.scatter(logistic\_data[}\StringTok{\textquotesingle{}time\textquotesingle{}}\NormalTok{], logistic\_data[}\StringTok{\textquotesingle{}y\textquotesingle{}}\NormalTok{], c}\OperatorTok{=}\StringTok{\textquotesingle{}black\textquotesingle{}}\NormalTok{, s}\OperatorTok{=}\DecValTok{10}\NormalTok{, zorder}\OperatorTok{=}\DecValTok{2}\NormalTok{)}
\NormalTok{plt.xticks((}\FloatTok{0.}\NormalTok{, }\FloatTok{2.5}\NormalTok{, }\FloatTok{5.}\NormalTok{, }\FloatTok{7.5}\NormalTok{, }\FloatTok{10.}\NormalTok{))}\OperatorTok{;}
\NormalTok{plt.yticks((}\FloatTok{0.}\NormalTok{, }\FloatTok{2.5}\NormalTok{, }\FloatTok{5.}\NormalTok{, }\FloatTok{7.5}\NormalTok{, }\FloatTok{10.}\NormalTok{))}\OperatorTok{;}
\NormalTok{plt.xlabel(}\StringTok{\textquotesingle{}Time\textquotesingle{}}\NormalTok{)}
\NormalTok{plt.ylabel(}\StringTok{\textquotesingle{}f\textquotesingle{}}\NormalTok{)}
\NormalTok{plt.tight\_layout()}
\end{Highlighting}
\end{Shaded}

\pandocbounded{\includegraphics[keepaspectratio]{01_fit_ode_files/figure-latex/unnamed-chunk-7-177.pdf}}

Here, the lines are the ODE, \(f\), that are solved using samples from
the posterior, \(p(\theta | y)\) and the dots show the data, \(y\).

We can also plot samples from the posterior predictive distribution,
\(p(y^{\rm rep} | y)\) (Gelman et al. 2013):

\[
p(y^{\rm rep} | y) = \int p(y^{\rm rep} | \theta)p(\theta | y) d\theta
\]

to check how well our entire model, with assumed priors and likelihood,
fits to the data. The fitted model should be able to generate data under
the posterior predictive distribution that resembles the observed data
(Gelman et al. 2013):

\section*{References}\label{references}
\addcontentsline{toc}{section}{References}

\protect\phantomsection\label{refs}
\begin{CSLReferences}{1}{1}
\bibitem[\citeproctext]{ref-gelman2013}
Gelman, Andrew, John B. Carlin, Hal S. Stern, David B. Dunson, Aki
Vehtari, and Donald B. Rubin. 2013. \emph{Bayesian Data Analysis, Third
Edition}. In \emph{Bayesian Data Analysis, Third Edition}. Chapman;
Hall/CRC.

\bibitem[\citeproctext]{ref-gelman2020}
Gelman, Andrew, Aki Vehtari, Daniel Simpson, et al. 2020. {``Bayesian
Workflow.''} \emph{arXiv Preprint arXiv:2011.01808}.

\bibitem[\citeproctext]{ref-Stan}
Stan Development Team. 2025. \emph{Stan Modeling Language Users Guide
and Reference Manual, Version 2.32}. \url{https://mc-stan.org}.

\bibitem[\citeproctext]{ref-vehtari2021}
Vehtari, Aki, Andrew Gelman, Daniel Simpson, Bob Carpenter, and
Paul-Christian Bürkner. 2021. {``Rank-Normalization, Folding, and
Localization: An Improved r ̂ for Assessing Convergence of MCMC (with
Discussion).''} \emph{Bayesian Analysis} 16 (2): 667--718.

\bibitem[\citeproctext]{ref-verhulst1838}
Verhulst, Pierre-François. 1838. {``Notice Sur La Loi Que La Population
Suit Dans Son Accroissement.''} \emph{Correspondence Mathematique Et
Physique} 10: 113--29.

\end{CSLReferences}

\end{document}
